\subsubsection{Detalhes Técnicos}

Os gráficos a seguir mostram a evolução da distribuição conforme o tamanho amostral aumenta.

\paragraph{Efeito da Assimetria na População Original}
A população Exponencial simulada possui assimetria fortemente positiva:

$\gamma_1 = 1.9684$

Este valor ($\gamma_1 > 2$) indica uma distribuição fortemente assimétrica à direita.

\paragraph{Validação do Teorema do Limite Central}
A assimetria das médias amostrais diminui com o aumento do tamanho amostral:

$$\gamma_1(n=5) = 0.8174$$
$$\gamma_1(n=50) = 0.3007$$

Para $n=50$, a assimetria é próxima de zero, comprovando que a distribuição das médias amostrais aproxima-se da Normal.

\paragraph{Validação da Cobertura do Intervalo de Confiança}
A taxa de cobertura do intervalo de confiança de 95\% é:

$$P(\mu \in [\overline{X} - 1.96 \cdot SE, \overline{X} + 1.96 \cdot SE]) = 0.9315$$

Este resultado (esperado $\approx 0.95$) confirma que o TCL permite o uso confiável dos intervalos de confiança baseados na distribuição Normal, mesmo quando a população original é assimétrica.

\paragraph{Conclusão}
O TCL é validado: a normalidade da distribuição das médias amostrais permite o uso confiável de métodos inferenciais baseados na distribuição Normal, independentemente da distribuição populacional.

